% 第 4 章
\bichapter{时钟树综合}{Clock Tree Synthesis}
\label{chap:cts}

要在 IC Compiler II 中执行时钟树综合（CTS），请参阅以下主题：
\begin{itemize}
  \item CTS 前提条件
  \item 定义时钟树
  \item 验证时钟树
  \item 设置时钟树设计规则约束
  \item 指定 CTS 设置
  \item 实现时钟树并执行 CTS 后优化
  \item 实现多源时钟树（Multisource Clock Trees）
  \item 分析时钟树结果
\end{itemize}

% ============================================================
\section{CTS 前提条件}
\subsection*{Prerequisites for Clock Tree Synthesis}

在对 block 运行 CTS 之前，应满足以下要求：
\begin{itemize}
  \item 已用 \cmd{create\_clock} 或 \cmd{create\_generated\_clock} 标识时钟源。
  \item block 已完成布局与优化。\\
        使用 \cmd{check\_legality -verbose} 确认布局合法。对布局不合法的 block 运行 CTS 可能导致运行时间过长、QoR 下降。\\
        开始 CTS 前，估计的 QoR 应已满足要求，包括：
        \begin{itemize}
          \item \textbf{拥塞（Congestion）}\\
                若 CTS 前拥塞未解决，插入时钟树可能使拥塞恶化。若 block 拥塞，可用 \opt{-congestion} 与 \opt{-congestion\_effort high} 重新运行 \cmd{create\_placement}，但运行时间可能较长。
          \item 时序（Timing）
          \item 最大电容（Maximum capacitance）
          \item 最大转换时间（Maximum transition time）
        \end{itemize}
  \item 电源与地网已预布线（prerouted）。
  \item 高扇出网络（如 scan enable）已用缓冲器综合。
  \item 已定义 active scenario。\\
        默认情况下，工具对所有为 setup 或 hold 分析启用的 active scenario 中的全部时钟进行综合与优化。
\end{itemize}

% ============================================================
\section{定义时钟树}
\subsection*{Defining the Clock Trees}

运行 CTS 前，应分析 block 中每棵时钟树并确定：
\begin{itemize}
  \item 时钟根（clock root）是什么；
  \item 所需的时钟 sink 与时钟树例外（exceptions）；
  \item 时钟树是否包含既有单元（如 clock-gating 单元）；
  \item 时钟树是否自收敛（\textbf{convergent clock path}）或与另一时钟树重叠（\textbf{overlapping clock path}）；
  \item 是否与其他时钟树存在时序关系（如 interclock skew 要求）。
\end{itemize}

据此定义时钟树并确认工具得到的定义正确。工具根据 block 中已定义的时钟推导时钟树；若不满足需求，可按下列主题定义：
\begin{itemize}
  \item 推导时钟树
  \item 定义时钟树例外
  \item 限制时钟网络上的优化
  \item 跨 mode 复制时钟树例外
  \item 由理想时钟延迟推导时钟树例外
  \item 处理存在平衡冲突的端点
\end{itemize}

\subsection{推导时钟树}
\subsubsection*{Deriving the Clock Trees}

工具从时钟根出发，沿传递扇出（transitive fanout）追踪至时钟端点以推导时钟树。一般在遇到时序单元或宏的时钟引脚时终止；但对集成时钟门控（ICG）单元，或其扇出驱动 generated clock 的时序单元，工具会继续穿越。

若门控逻辑使用非单极性（non-unate）单元（如 XOR/XNOR），追踪时会同时使用正单极性与负单极性时序弧。若与 block 功能模式不符，可用 \cmd{set\_case\_analysis} 将单元非时钟输入保持为常量，迫使其进入 CTS 所需功能模式。

例如，图~\ref{fig:cts-non-unate} 所示门控逻辑中，可用下列命令使 XOR 门（U0）进入功能模式：
\begin{lstlisting}
icc2_shell> set_case_analysis 0 U0/B
\end{lstlisting}

\begin{figure}[htbp]
\centering
\fbox{\parbox{0.85\textwidth}{\centering\vspace{1.2cm}
\textit{（原书 Figure 36：Non-Unate Gated Clock）}\\[0.3em]
\small 请参阅原版 PDF 插图；可将截图放入 \texttt{figures/} 后用 \texttt{\textbackslash includegraphics} 替换。
\vspace{1.2cm}}}
\caption{非单极性门控时钟}
\label{fig:cts-non-unate}
\end{figure}

\subsubsection{识别时钟根}
\subsubsection*{Identifying the Clock Roots}

工具将 \cmd{create\_clock} 定义的时钟源（输入端口或内部层次引脚）作为时钟根。

对由 \cmd{create\_generated\_clock} 定义的嵌套时钟树，工具将 master-clock 源视为时钟根，嵌套树端点视为该 master-clock 源的端点。例如图~\ref{fig:cts-nested} 中，工具将 CLK 端口视为 \cmd{genclk1} 的时钟根。

\begin{figure}[htbp]
\centering
\fbox{\parbox{0.85\textwidth}{\centering\vspace{1.2cm}
\textit{（原书 Figure 37：Nested Clock Tree With a Generated Clock）}
\vspace{1.2cm}}}
\caption{带 generated clock 的嵌套时钟树}
\label{fig:cts-nested}
\end{figure}

若 block 含 generated clock，须正确界定 master-clock 源，否则 skew 与时序 QoR 会变差：
\begin{itemize}
  \item 若无法回溯到 master-clock 源，则无法将该 generated clock 的 sink 与其源的 sink 一并平衡；
  \item 若 master-clock 源并非由 \cmd{create\_clock}/\cmd{create\_generated\_clock} 定义的时钟源，则无法为其或其源综合时钟树。
\end{itemize}
用 \cmd{check\_clock\_trees} 验证 master-clock 源定义，见「验证时钟树」。

\subsubsection{指定时钟根时序特性}
\subsubsection*{Specifying the Clock Root Timing Characteristics}

为获得真实 CTS 结果，须正确建模时钟根时序特性：
\begin{itemize}
  \item 若时钟根为无 I/O pad 的输入端口，须准确指定 driving cell。\\
        若驱动过弱，工具可能插入额外缓冲以满足最大转换时间、最大电容等 DRC。\\
        未指定 driving cell（或驱动强度）时，工具假定端口具有无限驱动能力。\\
        例如，CLK1 为 CLK1 时钟树根，设驱动为 mylib/CLKBUF：
\begin{lstlisting}
icc2_shell> set_driving_cell -lib_cell mylib/CLKBUF [get_ports CLK1]
\end{lstlisting}
  \item 若时钟根为已插入 I/O pad 的输入端口，须准确指定输入转换时间。例如上升 0.3、下降 0.2：
\begin{lstlisting}
icc2_shell> set_input_transition -rise 0.3 [get_ports CLK1]
icc2_shell> set_input_transition -fall 0.2 [get_ports CLK1]
\end{lstlisting}
\end{itemize}

\subsubsection{识别时钟端点}
\subsubsection*{Identifying the Clock Endpoints}

推导时钟树时，工具识别两类端点：
\begin{itemize}
  \item \textbf{Sink pins}（亦称 balancing pins）：用于延迟平衡。工具为所有 sink 赋予零插入延迟并用于平衡。CTS 中 sink 参与 DRC 与时钟时序（skew、insertion delay）的计算与优化。
  \item \textbf{Ignore pins}：排除在时钟时序计算与优化之外；仅参与 DRC 计算与优化。CTS 时在 ignore pin 前插入 guide buffer 以隔离；该点之外不再做 skew/insertion delay 优化，但仍做 DRC 修复。
\end{itemize}

工具将下列端点识别为 sink pins：
\begin{itemize}
  \item 时序单元（latch 或 flip-flop）的时钟引脚（除非该单元驱动 generated clock）；
  \item 宏单元的时钟引脚。
\end{itemize}

工具将下列端点识别为 ignore pins：
\begin{itemize}
  \item 位于另一时钟扇出中的时钟树源引脚。例如图~\ref{fig:cts-cascaded} 中，被驱动时钟 clk2 的源引脚是驱动时钟 clk1 的 ignore pin；被驱动时钟的 sink 不视为驱动时钟的 sink。
\begin{figure}[htbp]
\centering
\fbox{\parbox{0.85\textwidth}{\centering\vspace{1.0cm}
\textit{（原书 Figure 38：Cascaded Clock With Two Source Clocks）}
\vspace{1.0cm}}}
\caption{双源级联时钟}
\label{fig:cts-cascaded}
\end{figure}
  \item 时序单元的非时钟输入引脚；
  \item 三态使能引脚；
  \item 输出端口；
  \item 定义不正确的时钟引脚（例如无触发沿信息，或无到输出的时序弧）；
  \item 用 \cmd{set\_case\_analysis} 保持常量的缓冲器/反相器输入引脚；
\begin{noteBox}
若时钟源被 \cmd{set\_case\_analysis} 保持为常量，工具不会综合该时钟树。
\end{noteBox}
  \item 无扇出、或无已启用时序弧的组合单元/ICG 输入引脚。
\end{itemize}

可在 GUI 中检查 sink/ignore 识别是否正确（见「在 GUI 中分析时钟树」）。若默认正确，时钟树定义即完成；否则先排查禁用时序弧与 case analysis（\cmd{report\_disable\_timing}、\cmd{report\_case\_analysis}），删除导致错误定义的设置。必要时用 \cmd{set\_clock\_balance\_points} 覆盖默认 balancing/ignore 设置，见「定义时钟树例外」。

\subsection{定义时钟树例外}
\subsubsection*{Defining Clock Tree Exceptions}

时钟树例外是用户对工具为特定时钟推导的默认端点（balancing 与 ignore pins）所做的更改。工具支持：
\begin{itemize}
  \item \textbf{用户定义 sink pins}：在工具推导之外额外定义 sink（例如将工具推导的 ignore pin 定义为 sink）。见「定义 Sink Pins」。
  \item \textbf{用户定义插入延迟要求}：为 sink（推导或用户指定）指定特殊 insertion delay。见「定义插入延迟要求」。
  \item \textbf{用户定义 ignore pins}：排除工具推导为 sink 的端点（例如排除组合逻辑之后的所有分支）。见「定义 Ignore Pins」。
\end{itemize}

\subsubsection{定义 Sink Pins}
\subsubsection*{Defining Sink Pins}

语法：
\begin{lstlisting}
set_clock_balance_points
   -clock clock
   -consider_for_balancing true
   -balance_points pins
   [-corners corners]
\end{lstlisting}

\opt{-clock} 可选；省略则对所有时钟生效。用户指定 sink 的默认 insertion delay 为 0；覆盖方法见「定义插入延迟要求」。默认应用于所有 corner；用 \opt{-corners} 可限定。

例如将 U2/A 指定为当前 corner 下 CLK 的 sink：
\begin{lstlisting}
icc2_shell> set_clock_balance_points -clock [get_clocks CLK] \
   -consider_for_balancing true -balance_points [get_pins U2/A]
\end{lstlisting}

报告用 \cmd{report\_clock\_balance\_points}（同时报告用户定义的 sink 与 ignore）。移除用 \cmd{remove\_clock\_balance\_points}：
\begin{lstlisting}
icc2_shell> remove_clock_balance_points -clock [get_clocks CLK] \
   -balance_points [get_pins U2/A]
\end{lstlisting}

\subsubsection{定义插入延迟要求}
\subsubsection*{Defining Insertion Delay Requirements}

用 \cmd{set\_clock\_balance\_points} 的 \opt{-delay} 覆盖 sink 默认零相位延迟。工具将指定延迟（正或负）加到到达该 sink 的计算插入延迟上。语法：
\begin{lstlisting}
set_clock_balance_points
   [-clock clock]
   -delay delay
   [-rise] [-fall] [-early] [-late]
   -balance_points pins
   [-corners corners]
\end{lstlisting}

默认用于最长/最短路径以及上升/下降路径；用 \opt{-rise}/\opt{-fall}/\opt{-early}/\opt{-late} 控制适用范围。默认所有 corner；用 \opt{-corners} 限定。

例如对 U2/CLK 在上升最短路径使用 2.0~ns：
\begin{lstlisting}
icc2_shell> set_clock_balance_points -clock [get_clocks CLK] \
   -rise -early -delay 2.0 -balance_points [get_pins U2/CLK]
\end{lstlisting}

用 \cmd{report\_clock\_balance\_points} 报告用户定义插入延迟、带用户插入延迟的推导 sink，以及用户定义 ignore pins。

\subsubsection{定义 Ignore Pins}
\subsubsection*{Defining Ignore Pins}

语法：
\begin{lstlisting}
set_clock_balance_points
   [-clock clock]
   -consider_for_balancing false
   -balance_points pins
\end{lstlisting}

\opt{-clock} 可选。对图~\ref{fig:cts-ignore} 中 CLK 树，若要忽略 U2 之后的所有分支：
\begin{lstlisting}
icc2_shell> set_clock_balance_points -clock [get_clocks CLK] \
   -consider_for_balancing false -balance_points [get_pins U2/A]
\end{lstlisting}

\begin{figure}[htbp]
\centering
\fbox{\parbox{0.85\textwidth}{\centering\vspace{1.0cm}
\textit{（原书 Figure 39：User-Defined Ignore Pin）}
\vspace{1.0cm}}}
\caption{用户定义的 ignore pin}
\label{fig:cts-ignore}
\end{figure}

报告用 \cmd{report\_clock\_balance\_points}。CTS 时工具添加 guide buffer 隔离 ignore pin；后续数据通路优化可修复 guide buffer 之外的 DRC。移除：
\begin{lstlisting}
icc2_shell> remove_clock_balance_points -clock [get_clocks CLK] \
   -balance_points [get_pins U2/CLK]
\end{lstlisting}

\subsubsection{确保时钟树例外有效}
\subsubsection*{Ensuring Clock Tree Exceptions are Valid}

工具通过引脚属性 \cmd{is\_clock\_used\_as\_clock} 是否为 true 判断引脚是否在时钟网络上。若对属性为 false 的引脚用 \cmd{set\_clock\_balance\_points} 指定例外，工具会接受但在 CTS 中忽略。

考虑图~\ref{fig:cts-beyond-net} 的时钟网络：默认情况下，U2/A、U2/Y、FF3/D 的 \cmd{is\_clock\_used\_as\_clock} 为 false，U2/A 之后视为数据网络；U2/A 为 ignore pin，其扇出排除在 CTS 之外。CTS 会在 U2/A 前加 guide buffer；后续数据通路优化可修 DRC。

\begin{figure}[htbp]
\centering
\fbox{\parbox{0.85\textwidth}{\centering\vspace{1.0cm}
\textit{（原书 Figure 40：Specifying Exceptions on Pins Beyond the Clock Network）}
\vspace{1.0cm}}}
\caption{在时钟网络之外的引脚上指定例外}
\label{fig:cts-beyond-net}
\end{figure}

若希望 CTS 用时钟树约束一直修到 FF3/D 的 DRC：
\begin{enumerate}
  \item 将 FF3/D 显式设为 ignore pin：
\begin{lstlisting}
icc2_shell> set_clock_balance_points -clock [get_clocks CLK] \
    -consider_for_balancing false -balance_points [get_pins FF3/D]
\end{lstlisting}
  \item 用 \cmd{set\_sense -clock\_leaf} 强制将其视为时钟网络引脚：
\begin{lstlisting}
icc2_shell> set_sense -clock_leaf [get_pins FF3/D]
\end{lstlisting}
        该命令会：将扇入到该 clock leaf 的时钟分支上所有引脚的 \cmd{is\_clock\_used\_as\_clock} 设为 true；将该 leaf 扇出上所有引脚设为 false，且时钟不继续传播。\\
        检查属性：
\begin{lstlisting}
icc2_shell> get_attribute [get_pins FF3/D] is_clock_used_as_clock
\end{lstlisting}
\end{enumerate}

\subsection{限制时钟网络上的优化}
\subsubsection*{Restricting Optimization on the Clock Network}

可通过以下方式限制时钟网络部分区域的优化：
\begin{itemize}
  \item \textbf{Don't touch}：标识不做 CTS/优化的部分，可防止特定引脚之外被修改、特定 net 被缓冲、特定单元被改尺寸。见「设置 Don't Touch」。
  \item \textbf{Size-only}：仅允许单元改尺寸；除非 placement status 为 fixed，否则仍可搬移。见「设置 Size-Only」。
\end{itemize}

\subsubsection{设置 Don't Touch}
\subsubsection*{Setting Don't Touch Settings}

\begin{lstlisting}
icc2_shell> set_dont_touch_network -clock_only [get_pins pin_name]
\end{lstlisting}
\opt{-clock\_only} 仅作用于时钟网络；省略则同时作用于数据与时钟路径。清除：
\begin{lstlisting}
icc2_shell> set_dont_touch_network -clock_only [get_pins pin_name] -clear
\end{lstlisting}
也可对具体对象用 \cmd{set\_dont\_touch}：
\begin{lstlisting}
icc2_shell> set_dont_touch [get_cells cell_name] true
icc2_shell> set_dont_touch [get_nets -segments net_name] true
\end{lstlisting}
报告：\cmd{report\_dont\_touch -all}。
\begin{noteBox}
若单元 placement status 为 fixed，CTS 中按 don't touch 单元处理。
\end{noteBox}

\subsubsection{设置 Size-Only}
\subsubsection*{Setting Size-Only Settings}

\begin{lstlisting}
icc2_shell> set_size_only [get_cells cell_name] true
icc2_shell> report_size_only -all
\end{lstlisting}

\subsection{跨 Mode 复制时钟树例外}
\subsubsection*{Copying Clock Tree Exceptions Across Modes}

用 \cmd{set\_clock\_tree\_options} 的 \opt{-copy\_exceptions\_across\_modes}、\opt{-from\_mode}、\opt{-to\_mode}，指定将例外从某一 mode 复制到一个或多个等价 mode。例如将 M1 的例外复制到 M2、M3：
\begin{lstlisting}
icc2_shell> set_clock_tree_options -copy_exceptions_across_modes \
   -from_mode M1 -to_mode {M2 M3}
\end{lstlisting}

\subsection{由理想时钟延迟推导时钟树例外}
\subsubsection*{Deriving Clock Tree Exceptions From Ideal Clock Latencies}

若已对特定 sink 设置理想时钟延迟，可用 \cmd{derive\_clock\_balance\_points} 将其转换为时钟树例外：把 \cmd{set\_clock\_latency} 在 sink 上的理想延迟转为 \cmd{set\_clock\_balance\_points}。应在设计仍为理想时钟、CTS 之前运行。

默认行为：
\begin{itemize}
  \item 为所有启用 setup/hold（或两者）的 active scenario 中的所有理想 primary clock 生成 balance point；不为 generated clock 生成。可用 \opt{-clocks}/\opt{-corners} 限定。
  \item 计算每个 primary ideal clock 的参考延迟 $=$ \cmd{set\_clock\_latency} 在时钟上指定的理想 source 延迟 $+$ network 延迟。可用 \opt{-reference\_latency} 指定不同参考值。
  \item 将推导的约束应用到设计；用 \opt{-output} 可改为输出含 \cmd{set\_clock\_balance\_points} 的文件。
\end{itemize}

推导的 balance point 按时钟与 corner 区分。须确保 CTS 所用各 mode、至少 worst corner 上对所有时钟正确设置了 \cmd{set\_clock\_latency}。

公式：
\[
\text{sink 处 clock balance point 延迟} = \text{时钟参考延迟} - \text{sink 处总理想时钟延迟}
\]
其中总理想延迟 $=$ 时钟 source latency $+$ sink 的理想 network latency。

示例约束：
\begin{lstlisting}
icc2_shell> set_clock_latency -source 0.5 [get_clocks clk]
icc2_shell> set_clock_latency 1.0 [get_clocks clk]
icc2_shell> set_clock_latency 1.5 [get_pins reg1/CK]
icc2_shell> set_clock_latency 0.5 [get_pins reg2/CK]
\end{lstlisting}
运行 \cmd{derive\_clock\_balance\_points} 后得到：
\begin{lstlisting}
icc2_shell> set_clock_balance_points -delay -0.5 \
   -balance_points [get_pins reg1/CK] -clock clk -corners worst
icc2_shell> set_clock_balance_points -delay 0.5 \
   -balance_points [get_pins reg2/CK] -clock clk -corners worst
\end{lstlisting}
为避免冲突，不要对已用该命令推导 balance point 的时钟再手动施加 balance point。

\subsection{处理存在平衡冲突的端点}
\subsubsection*{Handling Endpoints With Balancing Conflicts}

某些 block 中存在结构上无法平衡的端点。工具可自动检测并推导 ignore pin 以解决冲突。对 \cmd{synthesize\_clock\_trees} 与 \cmd{clock\_opt} 默认启用；设 \appopt{cts.common.enable\_auto\_exceptions} 为 \cmd{false} 可关闭。

可检测并修复的冲突类型：
\begin{itemize}
  \item \textbf{含组合时钟路径模块的内部 sink}：若模块同时含内部 sink 与组合时钟路径，无法平衡内部 sink；工具将其定义为 ignore pin。例如图~\ref{fig:cts-comb-mod}(a) 网表对应 (b) 中带 check pin 的 ETM：ETM 内部 check pin 延迟来自 ETM 时序库，CTS 视其为 sink 并与顶层其他 sink 平衡；若到 check pin 为最短路径，顶层 CTS 无法在最短路径上插缓冲，导致大 skew。将其设为 ignore 可解决。
\begin{figure}[htbp]
\centering
\fbox{\parbox{0.85\textwidth}{\centering\vspace{1.0cm}
\textit{（原书 Figure 41：Module With Combinational Clock Path）}
\vspace{1.0cm}}}
\caption{含组合时钟路径的模块}
\label{fig:cts-comb-mod}
\end{figure}
  \item \textbf{无法同时平衡的 sink}：两时钟驱动公共 sink，且其一因 generated clock 成为另一 sink 的 parent 时，无法同时将 parent、child 与公共 sink 平衡。工具将 parent 时钟引脚定义为 ignore pin。例如图~\ref{fig:cts-unbal} 中，工具将 FF\_gen/CLK 定义为 clkb 的 ignore pin。
\begin{figure}[htbp]
\centering
\fbox{\parbox{0.85\textwidth}{\centering\vspace{1.0cm}
\textit{（原书 Figure 42：Clock Pins That Cannot Be Balanced Simultaneously）}
\vspace{1.0cm}}}
\caption{无法同时平衡的时钟引脚}
\label{fig:cts-unbal}
\end{figure}
\end{itemize}

% ============================================================
\section{验证时钟树}
\subsection*{Verifying the Clock Trees}

综合前用 \cmd{check\_clock\_trees} 验证定义。例如：
\begin{lstlisting}
icc2_shell> check_clock_trees -clocks [get_clocks CLK]
\end{lstlisting}
未指定 \opt{-clock} 时检查当前 block 全部时钟。检查项包括：
\begin{itemize}
  \item 无 sink 的时钟（master 或 generated）；
  \item 时钟网络成环；
  \item 重叠时钟使多时钟到达同一寄存器，但未启用 multiple-clocks-per-register 传播；
  \item 被忽略的时钟树例外；
  \item 在输出引脚上定义的 stop/float pin；
  \item 时钟树 reference 中使用了具有多条时序弧的缓冲器；
  \item 导致空缓冲列表的情形；
  \item generated clock 无有效 master-clock 源，包括：
        \begin{itemize}
          \item \cmd{create\_generated\_clock} 中指定的 master 不存在；
          \item master 未驱动 generated clock 的源引脚；
          \item 源引脚由多时钟驱动，但部分 master 未在 \cmd{create\_generated\_clock} 中指定。\\
                例如图~\ref{fig:cts-invalid-master}：GEN\_REG/Q 由 CLKA、CLKB 驱动，若仅指定 CLKA 为 master，则 GEN\_CLK 无有效 master。
        \end{itemize}
\end{itemize}

\begin{figure}[htbp]
\centering
\fbox{\parbox{0.85\textwidth}{\centering\vspace{1.0cm}
\textit{（原书 Figure 43：Generated Clock With Invalid Master Clock Source）}
\vspace{1.0cm}}}
\caption{master 源无效的 generated clock}
\label{fig:cts-invalid-master}
\end{figure}

对多角多模（MCMM）设计，还检查：
\begin{itemize}
  \item 冲突的 per-clock 例外设置；
  \item 冲突的 balancing 设置。
\end{itemize}
实现前应手动修复报告问题；每条消息有 man 页说明修复方法。修复全部问题可改善时钟树与时序 QoR。

% ============================================================
\section{设置时钟树设计规则约束}
\subsection*{Setting Clock Tree Design Rule Constraints}

CTS 支持以下设计规则约束：
\begin{itemize}
  \item \textbf{最大电容}：用 \cmd{set\_max\_capacitance} 的 \opt{-clock\_path} 指定。未指定时 CTS 默认为 0.6~pF。
  \item \textbf{最大转换时间}：用 \cmd{set\_max\_transition} 的 \opt{-clock\_path} 指定。未指定时 CTS 默认为 0.5~ns。
\end{itemize}

默认作用于当前 mode 关联的所有 corner。对特定 mode：用 \cmd{get\_clocks -mode}；对特定 corner：用约束命令的 \opt{-corners}。须跨 mode 及其 corner 正确施加。

例如对当前 mode 所有 corner 的 CLK 路径设最大转换 0.20~ns：
\begin{lstlisting}
icc2_shell> set_max_transition 0.20 -clock_path [get_clocks CLK]
\end{lstlisting}

对不同 corner 设不同值：
\begin{lstlisting}
icc2_shell> set_max_transition 0.15 -corners corner1 \
   -clock_path [get_clocks -mode mode1]
icc2_shell> set_max_transition 0.10 -corners corner2 \
   -clock_path [get_clocks -mode mode1]
\end{lstlisting}

对 mode 关联的不同 corner 设相同最大电容：
\begin{lstlisting}
icc2_shell> set_max_capacitance 0.6 \
    -clock_path [get_clocks -mode mode2] \
   -corners [get_corner \
      [get_attribute [get_modes mode2] associated_corners]]
\end{lstlisting}

% ============================================================
\section{指定 CTS 设置}
\subsection*{Specifying the Clock Tree Synthesis Settings}

CTS 选项指导时钟树实现，包括以下主题：
\begin{itemize}
  \item 指定时钟树 References
  \item 设置 Skew 与 Latency 目标
  \item 启用局部 Skew 优化
  \item 指定 CTS 主 Corner
  \item 阻止综合特定时钟
  \item 保留既有时钟树
  \item 启用时钟树降功耗技术
  \item 启用线长感知的功耗驱动搬移
  \item 降低电迁移（EM）
  \item 处理不准确约束
  \item 定义时钟单元间距规则
  \item 创建 Skew Groups
  \item 定义时钟单元名前缀
  \item 初始 CTS 使用全局路由器
  \item 为时钟网络指定约束
  \item 为特定端点构建更快时钟树
  \item 改进延迟驱动的时钟门分裂
  \item 降低时钟网络上的信号完整性影响
  \item 指定时钟延迟调整设置
  \item 报告时钟树设置
\end{itemize}

\subsection{指定时钟树 References}
\subsubsection*{Specifying the Clock Tree References}

可用于构建时钟树的缓冲器/反相器，以及既有门的 reference cell，称为时钟树 references。用 \cmd{set\_lib\_cell\_purpose -include cts} 指定。列表须满足：
\begin{itemize}
  \item 至少含一个缓冲器或一个反相器；
  \item 包含既有门的库单元（否则无法在 CTS/优化中对其改尺寸）。可用 \cmd{derive\_clock\_cell\_references} 自动推导非缓冲/反相既有单元的等价单元，见「为既有门推导 References」；
  \item 多电压且需 always-on 缓冲时，列表须含 always-on 单元；
  \item 列表中的库单元不得有 \cmd{dont\_touch}（即便列入 reference 也不会被 CTS 使用）。
\end{itemize}

为确保仅使用指定单元、且这些单元不被 CTS 以外的优化使用：
\begin{lstlisting}
icc2_shell> set cts_cells list_of_cells
icc2_shell> set_lib_cell_purpose -exclude cts [get_lib_cells]
icc2_shell> set_lib_cell_purpose -include none [get_lib_cells $cts_cells]
icc2_shell> set_lib_cell_purpose -include cts [get_lib_cells $cts_cells]
\end{lstlisting}

\subsubsection{为既有门推导 References}
\subsubsection*{Deriving Clock Tree References for Preexisting Gates}

须将等价单元指定为 clock references，否则影响 QoR。\cmd{derive\_clock\_cell\_references} 会为非缓冲/反相的既有时钟单元在等价库单元上设置 \cmd{valid\_purpose}；不推导缓冲器与反相器的等价单元，须手动指定：
\begin{lstlisting}
icc2_shell> set_lib_cell_purpose -include cts \
   {tech_lib/clk_buf* tech_lib/clk_inv*}
icc2_shell> derive_clock_cell_references
icc2_shell> synthesize_clock_trees
\end{lstlisting}

用 \opt{-output} 可生成 Tcl 脚本而非直接应用：
\begin{lstlisting}
icc2_shell> derive_clock_cell_references -output cts_leq_cells.tcl
icc2_shell> set_lib_cell_purpose -include cts \
   {tech_lib/clk_buf* tech_lib/clk_inv*}
icc2_shell> source cts_leq_cells.tcl
icc2_shell> synthesize_clock_trees
\end{lstlisting}

\subsubsection{限制所用目标库}
\subsubsection*{Restricting the Target Libraries Used}

用 \cmd{set\_target\_library\_subset -clock} 可限制顶层或逻辑层次下级所用库；须将 \appopt{opt.common.enable\_target\_library\_subset\_opt} 设为 1。例如全局允许 HVT/LVT 的 buf，但 TSK\_BLK 仅用 LVT\_lib：
\begin{lstlisting}
icc2_shell> set_lib_cell_purpose -include cts \
   {HVT_lib/buf1 HVT_lib/buf2 LVT_lib/buf1 LVT_lib/buf2}
icc2_shell> set_target_library_subset -clock {LVT_lib} \
       -objects [TOP/TSK_BLK]
icc2_shell> set_app_options \
  -name opt.common.enable_target_library_subset_opt -value 1
\end{lstlisting}

\subsection{设置 Skew 与 Latency 目标}
\subsubsection*{Setting Skew and Latency Targets}

默认 CTS 追求所有时钟的最佳 skew 与 latency，对低频时钟可能带来面积、功耗与运行时间开销。
\begin{itemize}
  \item skew 目标：\cmd{set\_clock\_tree\_options -target\_skew}
  \item latency 目标：\cmd{set\_clock\_tree\_options -target\_latency}
\end{itemize}
默认作用于所有时钟与所有 corner；用 \opt{-clocks}/\opt{-corners} 限定。报告：\cmd{report\_clock\_tree\_options}。移除：\cmd{remove\_clock\_tree\_options}（\opt{-all} 或相应选项）。

\subsection{启用局部 Skew 优化}
\subsubsection*{Enabling Local Skew Optimization}

默认最小化全局 skew（最长与最短时钟路径之差）。若最长/最短路径寄存器之间无时序路径，优化全局 skew 未必改善 QoR；优化局部 skew（时序路径 launch/capture 寄存器间最坏 skew）可改善 QoR。局部 skew 优化针对所有违例的 setup/hold 路径。

对 \cmd{synthesize\_clock\_trees} 与 \cmd{clock\_opt} 的综合与优化阶段启用：
\begin{lstlisting}
icc2_shell> set_app_options -list {cts.compile.enable_local_skew true}
icc2_shell> set_app_options -list {cts.optimize.enable_local_skew true}
\end{lstlisting}
启用后默认由工具推导有助于 QoR 的 skew 目标，并忽略 \cmd{set\_clock\_tree\_options -target\_skew}。若要阻止自动推导：
\begin{lstlisting}
icc2_shell> set_app_options \
 -name cts.common.enable_auto_skew_target_for_local_skew -value false
\end{lstlisting}

\subsection{指定 CTS 主 Corner}
\subsubsection*{Specifying the Primary Corner for Clock Tree Synthesis}

初始 CTS 默认识别时钟延迟最差的 corner，并在该 corner 的所有 mode 中插入缓冲以平衡延迟。用 \appopt{cts.compile.primary\_corner} 指定主 corner。

\subsection{阻止综合特定时钟}
\subsubsection*{Preventing Specific Clocks From Being Synthesized}

默认综合所有 \cmd{create\_clock}/\cmd{create\_generated\_clock} 定义的时钟。用 \appopt{cts.common.skip\_cts\_clocks} 排除：
\begin{lstlisting}
icc2_shell> set_app_options \
   -name cts.common.skip_cts_clocks -value CK_B
\end{lstlisting}

\subsection{保留既有时钟树}
\subsubsection*{Preserving Preexisting Clock Trees}

默认 CTS 会移除所有既有时钟缓冲器与反相器。将 \appopt{cts.compile.remove\_existing\_clock\_trees} 设为 \cmd{false} 可阻止。要保留特定既有缓冲/反相器，对其施加 don't touch 或 size-only 例外。

\subsection{启用时钟树降功耗技术}
\subsubsection*{Enabling Clock Tree Power Reduction Techniques}

用 \appopt{cts.compile.power\_opt\_mode} 启用，见表~\ref{tab:cts-power-opt}。

\begin{table}[htbp]
\centering
\caption{\appopt{cts.compile.power\_opt\_mode} 设置}
\label{tab:cts-power-opt}
\begin{tabularx}{\textwidth}{@{}l X@{}}
\toprule
\textbf{目的} & \textbf{设置} \\
\midrule
将 clock-gating 单元移近驱动以缩短高翻转输入网线长 & \cmd{gate\_relocation} \\
用聚类技术改善功耗 & \cmd{low\_power\_targets} \\
同时使用门控搬移与聚类 & \cmd{all} \\
\bottomrule
\end{tabularx}
\end{table}

\subsection{启用线长感知的功耗驱动搬移}
\subsubsection*{Enabling Wire Length Aware Power-Driven Relocation}

工具支持改进的活动驱动时钟门搬移：控制线长劣化，并在搬移前考虑功耗影响。将 \appopt{cts.compile.power\_opt\_mode} 设为 \cmd{all} 或 \cmd{gate\_relocation}。下列情况不搬移：
\begin{itemize}
  \item 搬移 ICG 功耗收益很小或恶化功耗，同时恶化线长；
  \item 搬移线长影响大而功耗收益小。
\end{itemize}

\subsection{降低电迁移}
\subsubsection*{Reducing Electromigration}

时钟单元功耗较高；局部聚集会提高电源轨电流密度，增加 EM 风险。可通过设置时钟单元间距，避免沿标准单元电源轨在垂直 strap 之间局部堆积：
\begin{enumerate}
  \item 用 \cmd{set\_clock\_cell\_spacing} 定义反相器、缓冲器、ICG 的间距规则（见「定义时钟单元间距规则」）；
  \item 用 \cmd{synthesize\_clock\_trees} 或 \cmd{clock\_opt} 执行 CTS；
\begin{noteBox}
\cmd{balance\_clock\_groups} 也遵守时钟单元间距规则。见「平衡不同时钟树之间的 Skew」。
\end{noteBox}
  \item 用 \cmd{check\_legality -verbose} 检查间距违例；设置合适约束后应无违例。
\end{enumerate}

\subsection{处理不准确约束}
\subsubsection*{Handling Inaccurate Constraints During Clock Tree Synthesis}

实现早期约束可能不准确或过严，直接 CTS 可能导致长运行时间或差 QoR。可设：
\begin{lstlisting}
icc2_shell> set_app_options \
   -name cts.common.enable_dirty_design_mode -value true
\end{lstlisting}
此时工具会：
\begin{itemize}
  \item 忽略时钟网上的 \cmd{dont\_touch}/\cmd{dont\_touch\_network}；
  \item 移除过紧的最大转换/最大电容约束；
  \item 移除小于 50~microns 的最大网长约束；
  \item 若水平间距 $>$ 3$\times$ site height 或垂直间距 $>$ 20$\times$ site height，忽略时钟单元间距规则；
  \item 若 NDR 的 width+spacing $>$ 10$\times$（默认 width+spacing），忽略 NDR；
  \item 报告大到足以增加时钟插入延迟的 balance point 延迟值。
\end{itemize}

\subsection{定义时钟单元间距规则}
\subsubsection*{Defining Clock Cell Spacing Rules}

用 \cmd{set\_clock\_cell\_spacing}；至少指定 x 或 y 方向最小间距（通常两者都指定）。
\begin{itemize}
  \item \opt{-x\_spacing}：微米，通常依赖驱动强度与时钟频率；
  \item \opt{-y\_spacing}：通常小于半行高。
\end{itemize}
\begin{noteBox}
间距过大因单元远离目标 sink 簇而增加 skew。
\end{noteBox}
默认作用于所有时钟单元；用 \opt{-lib\_cells}/\opt{-clocks} 限定。相邻两单元若同方向都有规则，间距为两者之和。可将 \appopt{cts.placement.cell\_spacing\_rule\_style} 设为 \cmd{maximum}，改为取较大值而非求和。

\cmd{synthesize\_clock\_trees}、\cmd{balance\_clock\_groups}、\cmd{clock\_opt} 均遵守这些规则。报告：\cmd{report\_clock\_cell\_spacings}。移除：\cmd{remove\_clock\_cell\_spacings}（默认全部；\opt{-lib\_cells} 可限定）。

\subsection{创建 Skew Groups}
\subsubsection*{Creating Skew Groups}

若希望一组 sink 仅在组内平衡，用 \cmd{create\_clock\_skew\_group}，以 \opt{-objects} 指定 sink。可选：\opt{-clock}、\opt{-mode}（默认当前 mode）、\opt{-name}。例如：
\begin{lstlisting}
icc2_shell> create_clock_skew_group -name sg1 \
   -objects {reg1/CP reg2/CP reg3/CP}
\end{lstlisting}
报告：\cmd{report\_clock\_skew\_groups}。移除：\cmd{remove\_clock\_skew\_groups}。

\subsection{定义时钟单元名前缀}
\subsubsection*{Defining a Name Prefix for Clock Cells}

用 \appopt{cts.common.user\_instance\_name\_prefix}：
\begin{lstlisting}
icc2_shell> set_app_options \
   -name cts.common.user_instance_name_prefix -value "CTS_"
\end{lstlisting}
数据网上优化（含 \cmd{clock\_opt}）的单元前缀用 \appopt{opt.common.user\_instance\_name\_prefix}。

\subsection{初始 CTS 使用全局路由器}
\subsubsection*{Using the Global Router During Initial Clock Tree Synthesis}

默认初始 CTS 用虚拟路由器；复杂 floorplan 可能产生拥塞热点并在时钟布线后劣化 QoR。将 \appopt{cts.compile.global\_route\_aware\_buffering} 设为 \cmd{true}，可在 \cmd{synthesize\_clock\_trees}/\cmd{clock\_opt} 的初始 CTS 阶段使用全局路由器。

\subsection{为时钟网络指定约束}
\subsubsection*{Specifying Constraints for Clock Nets}

默认用默认布线规则与可用布线层。可按「指定布线资源」设置时钟网最小/最大布线层，用于 RC 估计、拥塞分析与布线。为降低线延迟，可用宽线与高层金属（NDR）；须先定义规则并赋给时钟网（见「使用非默认布线规则」）。

\subsection{为特定端点构建更快时钟树}
\subsubsection*{Constructing Faster Clock Tree to Specific Endpoints}

可为指定端点构建无额外绕线/延迟平衡的更快时钟树以改善 latency。须先为所需端点创建用户 skew group，再运行 \cmd{synthesize\_clock\_trees} 或 \cmd{clock\_opt} 的 \cmd{build\_clock} 阶段：
\begin{lstlisting}
icc2_shell> create_clock_skew_group -objects <endpoints>
\end{lstlisting}
\begin{noteBox}
\begin{itemize}
  \item 若快速端点之间有平衡要求，须放入同一 skew group；
  \item 若无平衡要求，须为每个端点单独建 skew group；
  \item 所需 skew group 端点的上游时钟逻辑与默认组或其他用户 skew group 共享，可能导致为平衡其他组端点而绕线。
\end{itemize}
\end{noteBox}

\begin{figure}[htbp]
\centering
\fbox{\parbox{0.85\textwidth}{\centering\vspace{1.0cm}
\textit{（原书 Figure 44：Faster Clock Tree to Specific Endpoints）}
\vspace{1.0cm}}}
\caption{到特定端点的更快时钟树}
\label{fig:cts-faster}
\end{figure}

若一门同时驱动 skew group 与非 skew group 端点，skew group 端点可能因非组端点的平衡要求而变慢。除建 skew group 外，还须用 \cmd{split\_clock\_cells -to} 沿路径克隆门，以从时钟源到目标端点建立独立路径（在 CTS 之前）：
\begin{lstlisting}
icc2_shell> split_clock_cells -to $endpoints
\end{lstlisting}
也可用 \cmd{split\_clock\_cells -to <\$sinks> -latency\_driven} 自动识别最优克隆方案；端点须同时用 \cmd{create\_clock\_skew\_group} 建组。

\subsection{改进延迟驱动的时钟门分裂}
\subsubsection*{Improving Latency-Driven Clock-Gate Splitting}

延迟驱动的 CG 分裂现可自动识别并评估不同克隆搬移，以改善插入延迟并满足最大级数约束。启用：
将 \appopt{cts.icg.latency\_driven\_cloning} 设为 \cmd{true}。

\begin{figure}[htbp]
\centering
\fbox{\parbox{0.85\textwidth}{\centering\vspace{1.0cm}
\textit{（原书 Figure 45：Latency Driven Clock-Gate Splitting）}
\vspace{1.0cm}}}
\caption{延迟驱动的时钟门分裂}
\label{fig:cts-cg-split}
\end{figure}

\subsection{降低时钟网络上的信号完整性影响}
\subsubsection*{Reducing Signal Integrity Effects on Clock Nets}

可通过推导更大间距的 NDR 并赋给无拥塞区域的时钟网改善时序。启用：
\begin{lstlisting}
icc2_shell> set_app_options \
   -name cts.optimize.enable_congestion_aware_ndr_promotion -value true
\end{lstlisting}
新 NDR 间距取决于原规则：
\begin{itemize}
  \item 原 NDR 间距 $\leq$ 2$\times$ 默认间距：加倍，名称加后缀 \texttt{\_ext\_spacing}；
  \item 原 NDR 间距 $>$ 2$\times$ 默认：增加一个默认间距，后缀同上；
  \item 默认规则：间距加倍，新名为 \cmd{default\_rule\_equivalent\_ndr\_double\_spacing}。
\end{itemize}

\subsection{指定时钟延迟调整设置}
\subsubsection*{Specifying Settings for Clock Latency Adjustments}

I/O 端口通常由真实或虚拟时钟约束。CTS 后须更新约束 I/O 的时钟延迟以保持边界约束准确。用 \cmd{clock\_opt} 时工具自动更新；可用 \cmd{set\_latency\_adjustment\_options} 控制。例如按 PC\_CLK 更新 VCLK1/VCLK2，排除 VCLK3：
\begin{lstlisting}
icc2_shell> set_latency_adjustment_options \
   -reference_clock PC_CLK -clocks_to_update {VCLK1 VCLK2} \
   -exclude_clocks VCLK3
\end{lstlisting}
用 \cmd{synthesize\_clock\_trees} 时不会自动更新；CTS 后运行 \cmd{compute\_clock\_latency}（遵守上述设置）。

\subsection{报告时钟树设置}
\subsubsection*{Reporting the Clock Tree Settings}

用 \cmd{report\_clock\_settings}。默认报告当前 block 所有时钟的配置（DRC、目标 skew/latency）、references、单元间距规则、NDR。用 \opt{-clock} 限定时钟；用 \opt{-type}：
\begin{itemize}
  \item \cmd{configurations}：DRC、目标 skew/latency；
  \item \cmd{references}：时钟树 references；
  \item \cmd{spacing\_rules}：单元间距；
  \item \cmd{routing\_rules}：时钟网 NDR。
\end{itemize}



% ============================================================
\section{实现时钟树并执行 CTS 后优化}
\subsection*{Implementing Clock Trees and Performing Post-CTS Optimization}

执行 CTS 前应保存 block，以便必要时以同一起点 refinement 目标后重跑。可用方法：
\begin{itemize}
  \item 独立 CTS（\cmd{synthesize\_clock\_trees}）
  \item 用 \cmd{clock\_opt} 综合、优化并布线时钟树
  \item 控制并发时钟与数据（CCD）优化
  \item 分裂时钟单元
  \item 平衡不同时钟树之间的 Skew
  \item 使用机器学习数据的基于全局布线优化
  \item 布线时钟树
  \item 在 CTS、优化与时钟布线中插入 Via Ladder
  \item CTS 后将时钟标记为 propagated
  \item 执行布线后时钟树优化
  \item 执行电压优化
  \item 将时钟树标记为已综合
  \item 移除时钟树
\end{itemize}

\subsection{独立时钟树综合}
\subsubsection*{Performing Standalone Clock Trees Synthesis}

\begin{lstlisting}
icc2_shell> synthesize_clock_trees
\end{lstlisting}
该命令执行：
\begin{enumerate}
  \item \textbf{CTS}：此前对时钟网做虚拟布线并用 RC 估计确定时序；
  \item \textbf{时钟树优化}：此前对时钟网做全局布线并用 RC 提取；优化中对被修改的时钟网做增量全局布线。
\end{enumerate}
\begin{noteBox}
该命令不对时钟树做详细布线。
\end{noteBox}
默认作用于所有为 setup/hold 启用的 active scenario 中的全部时钟。用 \opt{-clock} 限定，例如：
\begin{lstlisting}
icc2_shell> synthesize_clock_trees -clocks [get_clocks CLK]
\end{lstlisting}
CTS 后，工具将所有 active scenario 各 mode 中的时钟设为 propagated。

\subsection{用 clock\_opt 综合、优化并布线时钟树}
\subsubsection*{Synthesizing, Optimizing, and Routing Clock Trees With the clock\_opt Command}

\cmd{clock\_opt} 用一条命令完成综合、布线与进一步优化，阶段为：
\begin{enumerate}
  \item \cmd{build\_clock}：综合并优化所有 active scenario 各 mode 的全部时钟；之后设为 propagated；
  \item \cmd{route\_clock}：对已综合时钟网做详细布线；
  \item \cmd{final\_opto}：进一步优化、时序驱动布局与合法化，然后全局布线并做广泛的基于全局布线的优化（含增量合法化与 route patching）。
\end{enumerate}
用 \opt{-from}/\opt{-to} 限制连续阶段；省略 \opt{-from} 从 \cmd{build\_clock} 开始，省略 \opt{-to} 跑到 \cmd{final\_opto} 完成。例如仅综合优化并详细布线：
\begin{lstlisting}
icc2_shell> clock_opt -to route_clock
\end{lstlisting}

因 \cmd{final\_opto} 做全局布线，运行前应设好所有路由器相关设置（布线规则、route guide、工艺节点应用选项等）。每个 block 应只跑一次 \cmd{final\_opto}；完成后信号网均已全局布线，后续 \cmd{route\_auto}/\cmd{route\_global}/\cmd{route\_group} 会跳过全局布线。为避免后续 track assignment/详细布线出错，\cmd{clock\_opt} 后不要改动任何全局布线形状。

\subsubsection{CTS 中考虑压降信息}
\subsubsection*{Considering Voltage Drop Information During Clock Tree Synthesis}

默认 \cmd{clock\_opt} 插入时钟缓冲时不考虑压降，可能把缓冲放在高 PG 电阻处。可用 RedHawk Fusion 分析 PG 网络：
\begin{enumerate}
  \item \cmd{analyze\_rail -min\_path\_resistance}；
  \item \cmd{open\_rail\_result} 加载结果；
  \item 将 \appopt{clock\_opt.flow.enable\_voltage\_drop\_aware} 设为 \cmd{true}；
  \item 运行 \cmd{clock\_opt}。
\end{enumerate}

\subsubsection{优化中对关键网络使用 NDR}
\subsubsection*{Using Nondefault Routing Rules for Critical Nets During Optimization}

布线前优化中可对时序关键网使用 NDR，并以软约束引导路由器。对 \cmd{clock\_opt} 启用：
\begin{lstlisting}
icc2_shell> set_app_options -name clock_opt.flow.optimize_ndr \
    -value true
\end{lstlisting}

\subsubsection{在 clock\_opt 中执行 CCD 优化}
\subsubsection*{Performing Concurrent Clock and Data Optimization During the clock\_opt Command}

在数据通路优化中利用正松弛调整寄存器时钟到达时间以改善时序 QoR，称为并发时钟与数据（CCD）优化。\cmd{clock\_opt} 默认启用；将 \appopt{clock\_opt.flow.enable\_ccd} 设为 \cmd{false} 可关闭。可选控制步骤：限制延迟调整值、排除边界路径、排除特定 path group/scenario/sink、控制时序/hold 优化努力、控制 I/O 时钟延迟调整、用 \cmd{report\_ccd\_timing} 分析效果（详见「控制 CCD 优化」各小节）。

\subsubsection{控制 clock\_opt 中的多比特优化}
\subsubsection*{Controlling Multibit Optimization Performed During the clock\_opt Command}

\begin{itemize}
  \item \textbf{混合驱动强度多比特单元重连}：违例路径经过较弱驱动位时，可重连使路径走较强驱动位。启用：\appopt{clock\_opt.flow.enable\_multibit\_rewiring}=\cmd{true}。
  \item \textbf{CTS 后多比特 debanking}：在 \cmd{final\_opto} 阶段，若不会引入 hold 违例则可 debank 以改善 QoR。启用：\appopt{clock\_opt.flow.enable\_multibit\_debanking}=\cmd{true}。
\end{itemize}

\subsubsection{在时钟网络上执行功耗或面积回收}
\subsubsection*{Performing Power or Area Recovery on the Clock Network}

启用 CCD 时，工具在 \cmd{final\_opto} 对时钟单元与寄存器做时钟功耗回收。未启用 CCD 时可：
\begin{lstlisting}
icc2_shell> set_app_options \
   -name clock_opt.flow.enable_clock_power_recovery \
   -value power
\end{lstlisting}
此前须：用 \cmd{set\_scenario\_status} 的 \opt{-dynamic\_power}/\opt{-leakage\_power} 启用功耗优化 scenario；（可选）用 \cmd{read\_saif} 提供翻转活动，否则用默认活动做动态功耗回收。面积回收则将值设为 \cmd{area}。

\subsubsection{在 clock\_opt 中执行 IR-Drop 感知布局}
\subsubsection*{Performing IR-Drop-Aware Placement During the clock\_opt Command}

\begin{enumerate}
  \item \cmd{clock\_opt -to route\_clock}；
  \item 配置 RedHawk Fusion 并用 \cmd{analyze\_rail -voltage\_drop} 做静态/动态压降分析；
  \item 将 \appopt{place.coarse.ir\_drop\_aware} 设为 \cmd{true}；
  \item （可选）按「控制 IR-Drop 感知布局」做附加设置；
  \item \cmd{clock\_opt -from final\_opto}。
\end{enumerate}

\subsection{控制并发时钟与数据优化}
\subsubsection*{Controlling Concurrent Clock and Data Optimization}

CCD 对 \cmd{place\_opt} 与 \cmd{clock\_opt} 默认启用；对 \cmd{route\_opt} 须将 \appopt{route\_opt.flow.enable\_ccd} 设为 \cmd{true}。可控制：
\begin{itemize}
  \item 限制延迟调整值
  \item 排除边界路径 / 特定 path group / scenario / sink
  \item 控制时序与 hold 优化努力
  \item 控制 I/O 时钟延迟调整
  \item 在 \cmd{route\_opt} 中做动态压降驱动的 CCD
  \item 在布线前/布线后阶段指定优化目标
  \item 布线后启用缓冲移除
  \item 报告 CCD 时序；将 CCD offset 缩放到新 scenario
  \item 对 latch 与离散时钟门一并 skew
\end{itemize}

\subsubsection{限制延迟调整值}
\subsubsection*{Limiting the Latency Adjustment Values}

\begin{itemize}
  \item 提前量上限：\appopt{ccd.max\_prepone}
  \item 推迟量上限：\appopt{ccd.max\_postpone}
\end{itemize}
无默认值，单位为库时序单位。例如：
\begin{lstlisting}
icc2_shell> set_app_options -list {ccd.max_prepone 0.2}
icc2_shell> set_app_options -list {ccd.max_postpone 0.1}
\end{lstlisting}

\subsubsection{排除边界路径}
\subsubsection*{Excluding Boundary Paths}

默认对 block 内所有路径做 CCD。边界寄存器为输入端口传递扇出、输出端口传递扇入中的寄存器。将 \appopt{ccd.optimize\_boundary\_timing} 设为 \cmd{false} 可排除边界路径。可用 \appopt{ccd.ignore\_ports\_for\_boundary\_identification} 使部分端口的边界路径仍参与 CCD：
\begin{lstlisting}
icc2_shell> set_app_options -name ccd.optimize_boundary_timing \
   -value false
icc2_shell> set_app_options \
   -name ccd.ignore_ports_for_boundary_identification \
   -value {IN_A OUT_A}
\end{lstlisting}
即便禁用边界寄存器 CCD，若改善共享时钟路径的内部寄存器 QoR，工具仍可优化其时钟树扇入锥。将 \appopt{ccd.optimize\_boundary\_timing\_upstream} 设为 \cmd{false} 可阻止，但会严重限制 CCD 范围。上述设置影响 \cmd{place\_opt}/\cmd{clock\_opt}/\cmd{route\_opt}。

\subsubsection{排除特定 Path Groups / Scenarios / Sinks}
\subsubsection*{Excluding Specific Path Groups, Scenarios, and Sinks}

\begin{lstlisting}
icc2_shell> set_app_options -name ccd.skip_path_groups \
   -value {CLK1 {CLK2 scnA}}
icc2_shell> set_app_options -name ccd.ignore_scenarios \
   -value {scnB}
\end{lstlisting}

排除特定 sink：
\begin{enumerate}
  \item 对 sink 设 \cmd{cts\_fixed\_balance\_pin}：
\begin{lstlisting}
icc2_shell> set_attribute -objects reg21/CK \
   -name cts_fixed_balance_pin -value true
\end{lstlisting}
  \item 将 \appopt{ccd.respect\_cts\_fixed\_balance\_pins} 设为 \cmd{true}（阻止调整这些 sink 的延迟；其他 sink 的 CCD 仍可能间接改变其延迟）。设为 \cmd{upstream} 可阻止从这些 sink 到时钟根的时钟路径被改动。
\end{enumerate}

\subsubsection{控制时序与 Hold 优化努力}
\subsubsection*{Controlling Timing and Hold Optimization Effort}

\begin{itemize}
  \item \appopt{ccd.timing\_effort}：\cmd{low}/\cmd{high}，默认 \cmd{medium}（影响 \cmd{clock\_opt} 的 \cmd{final\_opto} 与 \cmd{route\_opt}）；
  \item \appopt{ccd.hold\_control\_effort}：\cmd{medium}/\cmd{high}/\cmd{ultra}，默认 \cmd{low}。提高 hold 优先级会减少可修复的 setup 违例，仅在 hold 关键时使用。
\end{itemize}

\subsubsection{控制 I/O 时钟延迟调整}
\subsubsection*{Controlling the Adjustment of I/O Clock Latencies}

\cmd{clock\_opt} 做 CCD 时默认调整约束边界端口的 I/O 时钟延迟；随后运行 \cmd{compute\_clock\_latency} 时也使用 CCD 技术。不更新的情形包括：\appopt{ccd.adjust\_io\_clock\_latency}=\cmd{false}、禁用边界路径 CCD、用 \appopt{ccd.skip\_path\_groups} 排除对应 path group、用 \cmd{set\_latency\_adjustment\_options -exclude\_clocks} 排除特定 I/O 时钟。

\subsubsection{route\_opt 中动态压降驱动的 CCD}
\subsubsection*{Performing Dynamic-Voltage-Drop-Driven CCD During route\_opt}

将 \appopt{route\_opt.flow.enable\_voltage\_drop\_opt\_ccd} 设为 \cmd{true}。工具用 RedHawk Fusion 识别动态压降热点与可搬移的 cell row，在优化中搬移或缩小尺寸以降低峰值动态压降且不伤害时序 QoR。可选阈值：
\begin{itemize}
  \item \appopt{ccd.voltage\_drop\_voltage\_threshold}：相对供电电压的压降百分比；
  \item \appopt{ccd.voltage\_drop\_population\_threshold}：按最差违例选取的单元数百分比。
\end{itemize}
另可将 \appopt{route\_opt.flow.enable\_irdrivenopt} 与上述选项同时设为 \cmd{true}，启用 IR 驱动 sizing（用更小漏电单元替换）。

\subsubsection{布线前阶段指定优化目标}
\subsubsection*{Specifying Optimization Targets at the Preroute Stage}

对 \cmd{place\_opt}/\cmd{clock\_opt} 的 CCD，可提高下列目标的 WNS 优先级：
\begin{itemize}
  \item 特定 path group：\appopt{ccd.targeted\_ccd\_path\_groups}（若同时在 \appopt{ccd.skip\_path\_groups} 中则仍跳过）；
  \item 端点文件：\appopt{ccd.targeted\_ccd\_end\_points\_file}（同时指定 path group 与端点时，仅优化指定组内的指定端点）；
  \item 最差 300 条路径：\appopt{ccd.enable\_top\_wns\_optimization}=\cmd{true}。
\end{itemize}

\subsubsection{布线后阶段指定优化目标}
\subsubsection*{Specifying Optimization Targets at the Postroute Stage}

\begin{enumerate}
  \item 将 \appopt{route\_opt.flow.enable\_targeted\_ccd\_wns\_optimization} 设为 \cmd{true}；
  \item 用 \appopt{ccd.targeted\_ccd\_path\_groups} 和/或 \appopt{ccd.targeted\_ccd\_end\_points\_file} 指定目标（至少其一）；
  \item （可选）\appopt{ccd.targeted\_ccd\_select\_optimization\_moves}：\cmd{auto}（默认，含 buffering）、\cmd{size\_only}、\cmd{equal\_or\_smaller\_sizing}、\cmd{footprint\_sizing}；
  \item （可选）\appopt{ccd.targeted\_ccd\_wns\_optimization\_effort}：\cmd{high}（默认）/\cmd{medium}/\cmd{low}；
  \item （可选）\appopt{ccd.targeted\_ccd\_threshold\_for\_nontargeted\_path\_groups}：非目标 path group 可被劣化的松弛阈值（默认正松弛可劣化至零）。
\end{enumerate}
可同时启用常规 CCD（\appopt{route\_opt.flow.enable\_ccd}）与 targeted CCD；此时先常规后 targeted。示例脚本（设计已布线并做过布线后优化）：
\begin{lstlisting}
open_lib design1
open_block blk.post_route
set_app_options -name route_opt.flow.enable_ccd -value false
set_app_options \
  -name route_opt.flow.enable_targeted_ccd_wns_optimization -value true
set_app_options -name ccd.targeted_ccd_path_groups -value clkA
set_app_options -name ccd.targeted_ccd_end_points_file -value clkA_ep.txt
set_app_options \
   -name ccd.targeted_ccd_select_optimization_moves -value size_only
set_app_options -name ccd.targeted_ccd_wns_optimization_effort \
   -value medium
route_opt
\end{lstlisting}

\subsubsection{布线后启用缓冲移除}
\subsubsection*{Enabling Buffer Removal at the Postroute Stage}

\begin{lstlisting}
icc2_shell> set_app_options -name route_opt.flow.enable_ccd -value true
icc2_shell> set_app_options -name ccd.post_route_buffer_removal -value true
icc2_shell> route_opt
\end{lstlisting}

\subsubsection{报告 CCD 时序}
\subsubsection*{Reporting Concurrent Clock and Data Timing}

用 \cmd{report\_ccd\_timing}。默认报告 block 中五个最关键端点寄存器的最差 capture（D-slack）与 launch（Q-slack）的 setup/hold。\opt{-type}：\cmd{fanin}/\cmd{fanout}、\cmd{stage}、\cmd{chain}（见图~\ref{fig:cts-ccd-stages}）。默认报告 setup；\opt{-hold} 报告 hold。\opt{-scenarios} 限定 scenario；\opt{-pins} 限定端点。可用 \opt{-prepone}/\opt{-postpone} 与 \opt{-annotate\_cell\_delay}/\opt{-annotate\_net\_delay} 做 what-if 分析（仅报告，不保存）。

\begin{figure}[htbp]
\centering
\fbox{\parbox{0.85\textwidth}{\centering\vspace{1.0cm}
\textit{（原书 Figure 46：Previous, Current, and Next Path Stages）}
\vspace{1.0cm}}}
\caption{前一、当前与下一路径阶段}
\label{fig:cts-ccd-stages}
\end{figure}

\subsubsection{将 CCD Offset 缩放到新 Scenario}
\subsubsection*{Scaling CCD Offsets to New Scenarios for Reporting Timing}

用 \cmd{copy\_useful\_skew} 将 CCD offset 复制到新 scenario，以获得一致的时序/QoR 报告。当 pin 的 CTS balance point offset 仅在部分 active scenario 中推导时，用该命令复制到无 offset 的 inactive/active scenario，并确保在计算新 offset 的所有 pin 与 scenario 上设置相应 \cmd{set\_clock\_latency}。先对有 balance point offset 的 scenario 取平均，再按各 scenario corner 的缩放因子缩放。

\subsubsection{对 Latch 与离散时钟门一并 Skew}
\subsubsection*{Skewing Latch and Discrete Clock Gates}

在 \cmd{compile\_fusion} 或 \cmd{place\_opt} 的 \cmd{final\_opto} 前将 \appopt{ccd.skew\_opt\_merge\_dcg\_cells\_by\_driver} 设为 \cmd{true}（默认 \cmd{false}），使 latch 与门控元件使用相同 offset，保持同级并改善离散时钟门的 hold。

\subsection{分裂时钟单元}
\subsubsection*{Splitting Clock Cells}

\begin{lstlisting}
icc2_shell> split_clock_cells -cells [get_cells U1/ICG*]
\end{lstlisting}
不分裂的情形：无 DRC 违例；有 don't-touch/size-only/fixed-placement；须在电源域或 \cmd{set\_freeze\_ports} 标识的 block 边界 punch port。分裂后新单元命名为 \texttt{<原名>\_split\_<整数>}，并复制原单元全部设置与约束。也可用 \opt{-loads} 指定由同一驱动驱动的负载集合：
\begin{lstlisting}
icc2_shell> set loads1 [get_pins I1/reg*/CK]
icc2_shell> set loads2 [get_pins I2/reg*/CK]
icc2_shell> split_clock_cells -loads [list $loads1 $loads2]
\end{lstlisting}

\subsection{平衡不同时钟树之间的 Skew}
\subsubsection*{Balancing Skew Between Different Clock Trees}

参与延迟平衡的时钟集合称为 clock balance group；可定义多个组，并为组内时钟定义延迟 offset；合称 interclock delay balancing 约束。
\begin{noteBox}
工具无法在 generated clock 与其他时钟之间平衡 skew。
\end{noteBox}
步骤：生成约束（手动或自动）→ 运行平衡。

\subsubsection{定义 Interclock 延迟平衡约束}
\subsubsection*{Defining the Interclock Delay Balancing Constraints}

用 \cmd{create\_clock\_balance\_group}：至少 \opt{-name} 与 \opt{-objects}（可用 \cmd{get\_clocks -mode}）。默认目标为零 offset，以组内最长插入延迟平衡。用 \opt{-offset\_latencies} 指定相对 offset；\opt{-corner} 限定 corner。例如 clk3 比 clk1/clk2 少 100：
\begin{lstlisting}
icc2_shell> create_clock_balance_group -name group3 \
   -objects [get_clocks {clk1 clk2 clk3}] \
   -offset_latencies {0 0 -100}
\end{lstlisting}
报告：\cmd{report\_clock\_balance\_groups}。移除：\cmd{remove\_clock\_balance\_groups}（指定组或 \opt{-all}）。

\subsubsection{自动生成约束}
\subsubsection*{Generating Interclock Delay Balancing Constraints Automatically}

\begin{lstlisting}
icc2_shell> derive_clock_balance_constraints
\end{lstlisting}
将有 interclock 时序路径的时钟放入同一 balance group。可用 \opt{-slack\_less\_than} 仅考虑松弛小于阈值的路径，例如：
\begin{lstlisting}
icc2_shell> derive_clock_balance_constraints -slack_less_than -0.2
\end{lstlisting}

\subsubsection{运行 Interclock 延迟平衡}
\subsubsection*{Running Interclock Delay Balancing}

\begin{lstlisting}
icc2_shell> balance_clock_groups
\end{lstlisting}
对 MCMM，对所有 active scenario 执行。

\subsection{使用机器学习数据的基于全局布线优化}
\subsubsection*{Performing Global-Route-Based Optimization Using Machine Learning Data}

对布线前/布线后时序相关性差的设计，可在 CTS 后用详细布线后得到的 ML 数据做优化。术语：Labels（预测输出）、Features（影响输入）、Model（映射关系）、Training（学习过程）。

\textbf{迭代 N（采集与训练）：}
\begin{enumerate}
  \item 完成 \cmd{clock\_opt} 的 \cmd{build\_clock} 与 \cmd{route\_clock}；
  \item 设 \appopt{est\_delay.ml\_delay\_opto\_dir} 为输出目录；
  \item 将 \appopt{est\_delay.ml\_delay\_gre\_mode} 设为 \cmd{feature}；
  \item \cmd{clock\_opt -from final\_opto}；
  \item 详细布线；
  \item 详细布线后、布线后优化前：\cmd{estimate\_delay -training\_labels}；
  \item \cmd{estimate\_delay -train\_model}；
  \item 继续布线后优化与后续流程。
\end{enumerate}

\textbf{迭代 N+1（应用模型）：} 同样完成时钟综合与布线 → 指定 ML 目录 → 将 \appopt{est\_delay.ml\_delay\_gre\_mode} 设为 \cmd{enable} → \cmd{clock\_opt -from final\_opto} → 详细布线 → \cmd{estimate\_delay -disable\_model} → 继续流程。

须保证各步骤/迭代中 active scenario 相同；设计、环境或流程变更时须重建模型。

\subsection{布线时钟树}
\subsubsection*{Routing Clock Trees}

\begin{lstlisting}
icc2_shell> route_group -all_clock_nets -reuse_existing_global_route true
\end{lstlisting}
\opt{-reuse\_existing\_global\_route true} 使详细路由器使用既有时钟全局布线，相关性更好。或：
\begin{lstlisting}
icc2_shell> clock_opt -from route_clock -to route_clock
\end{lstlisting}

\subsection{在 CTS、优化与时钟布线中插入 Via Ladder}
\subsubsection*{Inserting Via Ladders During CTS, Optimization, and Clock Routing}

Via ladder 是从引脚层堆叠到上层、供路由器连接的通孔叠层，可降低通孔电阻，改善性能与 EM 鲁棒性。流程：
\begin{enumerate}
  \item 定义 via ladder 规则；
  \item 用 \cmd{set\_via\_ladder\_candidate} 指定库引脚可用的 via ladder；
  \item 用 \cmd{set\_via\_ladder\_rules} 定义约束；
  \item 将 \appopt{opt.common.enable\_via\_ladder\_insertion} 设为 \cmd{true}；
  \item \cmd{clock\_opt -to build\_clock}；
  \item \cmd{insert\_via\_ladders}；
  \item \cmd{clock\_opt -from route\_clock}。
\end{enumerate}

\subsection{CTS 后将时钟标记为 Propagated}
\subsubsection*{Marking Clocks as Propagated After Clock Tree Synthesis}

运行 \cmd{synthesize\_clock\_trees} 或 \cmd{clock\_opt} 后，工具移除所有 active scenario 中已综合时钟的 ideal 设置，并将对应寄存器时钟引脚设为 propagated。若有 CTS 前未设为 active 的 scenario：
\begin{lstlisting}
icc2_shell> set_scenario_status -active true [all_scenarios]
icc2_shell> synthesize_clock_trees -propagate_only
\end{lstlisting}
\opt{-propagate\_only} 仅移除 ideal 设置，不执行 CTS。

\subsection{执行布线后时钟树优化}
\subsubsection*{Performing Postroute Clock Tree Optimization}

详细布线时钟/信号网后，因全局布线与详细布线差异以及耦合/串扰，时钟 QoR 可能下降。用 \cmd{synthesize\_clock\_trees -postroute}，并以 \opt{-routed\_clock\_stage} 指明布线阶段：
\begin{lstlisting}
icc2_shell> synthesize_clock_trees -postroute \
   -routed_clock_stage detail
icc2_shell> synthesize_clock_trees -postroute \
   -routed_clock_stage detail_with_signal_routes
\end{lstlisting}
\begin{noteBox}
启用 CCD（含用该技术的功耗/面积回收）时，布线后时钟树优化不优化 latency/skew，仅修复时钟网络上的逻辑 DRC。
\end{noteBox}

\subsection{执行电压优化}
\subsubsection*{Performing Voltage Optimization}

电压优化尝试在更低电压下获得更好的 PPA。需要：不同电压的 scaled library panes；以及带 scaled library 设置的 signoff 方法。高层流程：
\begin{enumerate}
  \item 原电压 \cmd{place\_opt}；
  \item 原电压 \cmd{clock\_opt}；
  \item \cmd{define\_scaling\_lib\_group} 关联相关 library panes；
  \item \cmd{set\_vopt\_range} 指定 corner 电压范围；
  \item \cmd{set\_vopt\_target} 指定 TNS/功耗目标与容差（可用 \opt{-excludeIO}）；
  \item \cmd{voltage\_opt}；
  \item 优化电压下 \cmd{route\_auto}；
  \item 优化电压下 \cmd{route\_opt}。
\end{enumerate}

\subsection{将时钟树标记为已综合}
\subsubsection*{Marking Clock Trees as Synthesized}

用 \cmd{mark\_clock\_trees} 防止工具修改既有时钟树，见表~\ref{tab:mark-cts}。工具遍历时钟树直至遇到例外引脚或默认 sink。

\begin{table}[htbp]
\centering
\caption{使用 \cmd{mark\_clock\_trees} 命令}
\label{tab:mark-cts}
\begin{tabularx}{\textwidth}{@{}l X@{}}
\toprule
\textbf{目的} & \textbf{选项} \\
\midrule
仅标记特定时钟（默认标记所有 active scenario 各 mode 中 \cmd{create\_clock}/\cmd{create\_generated\_clock} 定义的时钟） & \opt{-clocks} \\
标记为已综合（默认） & \opt{-synthesized} \\
清除 synthesized 属性 & \opt{-clear} \\
施加 \cmd{dont\_touch} & \opt{-dont\_touch} \\
清除 \cmd{dont\_touch} & \opt{-clear -dont\_touch} \\
传播 \cmd{set\_clock\_routing\_rules} 的 NDR & \opt{-routing\_rules} \\
传播 \cmd{set\_clock\_cell\_spacing} 间距规则 & \opt{-clock\_cell\_spacing} \\
将 sink 标记为 fixed & \opt{-fix\_sinks} \\
冻结时钟网布线 & \opt{-freeze\_routing} \\
\bottomrule
\end{tabularx}
\end{table}

\subsection{移除时钟树}
\subsubsection*{Removing Clock Trees}

\cmd{remove\_clock\_trees} 从根到 sink 遍历并移除缓冲器与反相器（保留带 \cmd{dont\_touch}/\cmd{size\_only} 的单元）。例外之外的单元仍视为时钟树一部分；clock-to-data 引脚之外视为数据通路不移除。同时重置 CTS 相关属性。
\begin{noteBox}
\cmd{remove\_clock\_trees} 不支持 clock mesh 网。
\end{noteBox}
默认移除所有 mode 的所有时钟树；用 \opt{-clocks} 限定（默认当前 mode；可用 \cmd{get\_clocks -mode}）。结构对移除的影响见表~\ref{tab:rm-cts}。默认同时移除被删时钟网的详细布线形状；可将网的 \cmd{shape\_use} 设为 \cmd{user\_route} 以保留。

\begin{table}[htbp]
\centering
\caption{时钟树移除行为}
\label{tab:rm-cts}
\begin{tabularx}{\textwidth}{@{}l X@{}}
\toprule
\textbf{对象} & \textbf{对移除的影响} \\
\midrule
Boundary cell & 移除 \\
Don't touch net 上的单元 & 保留 \\
Don't touch / Fixed cell & 保留 \\
Generated clock & 若定义在缓冲/反相器引脚上则保留；遍历与移除继续穿过 \\
Guide buffer & 移除 \\
ICG & 保留；遍历继续穿过 \\
Block abstraction model & 保留；遍历继续穿过 \\
Inverter & 仅成对移除；单独反相器不移除 \\
Isolation / Level shifter & 保留 \\
Three-state buffer & 保留；遍历在此停止 \\
例外之外添加的缓冲/反相器 & 移除 \\
\bottomrule
\end{tabularx}
\end{table}

% ============================================================
\section{实现多源时钟树}
\subsection*{Implementing Multisource Clock Trees}

本节介绍多源时钟树结构类型及实现流程：
\begin{itemize}
  \item 多源时钟树结构简介
  \item 实现常规多源时钟树
  \item 使用集成 Tap Assignment 的常规多源时钟树
  \item 仅含 H-Tree 全局结构的常规多源时钟树
  \item 实现 MSCTS 多重 H-Tree
  \item 实现功耗感知柔性 H-Tree
  \item 实现结构型多源时钟树（SMSCTS）
  \item 使用集成子树综合的 SMSCTS
  \item 改善 SMSCTS 与 CCD 的相关性
  \item 插入时钟驱动器；综合全局时钟树；创建时钟 strap；布线到 strap；分析 clock mesh
  \item 自动 Tap 插入与 H-Tree 综合；指定 Tap Assignment 选项；构建本地子树
\end{itemize}

\subsection{多源时钟树结构简介}
\subsubsection*{Introduction to Multisource Clock Trees Structures}

多源时钟树对片上偏差更宽容，跨 corner 性能优于传统结构，由以下组成：
\begin{itemize}
  \item \textbf{全局时钟结构}：时钟根、通常为 H-tree 的全局时钟树、mesh drivers、clock mesh；
  \item \textbf{本地子树}，驱动方式可为：
        \begin{itemize}
          \item 由连接到 mesh 的预定义 tap drivers 驱动（图~\ref{fig:mscts-reg}）——子树用常规 CTS（\cmd{synthesize\_clock\_trees}/\cmd{clock\_opt}）构建，称\textbf{常规多源时钟树}；
          \item 直接从 mesh 多点驱动（图~\ref{fig:mscts-str}）——保留用户结构并通过合并/分裂优化，称\textbf{结构型多源时钟树}。
        \end{itemize}
\end{itemize}

\begin{figure}[htbp]
\centering
\fbox{\parbox{0.85\textwidth}{\centering\vspace{1.0cm}
\textit{（原书 Figure 47：Regular Multisource Clock Tree）}
\vspace{1.0cm}}}
\caption{常规多源时钟树}
\label{fig:mscts-reg}
\end{figure}

\begin{figure}[htbp]
\centering
\fbox{\parbox{0.85\textwidth}{\centering\vspace{1.0cm}
\textit{（原书 Figure 48：Structural Multisource Clock Tree）}
\vspace{1.0cm}}}
\caption{结构型多源时钟树}
\label{fig:mscts-str}
\end{figure}

\subsection{实现常规多源时钟树}
\subsubsection*{Implementing a Regular Multisource Clock Tree}

\begin{enumerate}
  \item 指定时钟树约束与设置；
  \item 多电压设计：启用 \appopt{cts.multisource.enable\_full\_mv\_support}；启用 \appopt{opt.common.allow\_physical\_feedthrough}；运行 \cmd{check\_mv\_design} 并修复；
  \item 启用 \appopt{cts.multisource.cell\_em\_aware}；
  \item \cmd{create\_clock\_drivers} 插入 tap drivers；
  \item \cmd{create\_clock\_straps} 创建 clock mesh；
  \item \cmd{create\_clock\_drivers} 插入 mesh drivers；
  \item \cmd{synthesize\_multisource\_global\_clock\_trees} 构建全局树；
  \item \cmd{route\_clock\_straps} 连接 mesh 与 tap；
  \item \cmd{analyze\_subcircuit} 分析 mesh；
  \item 指定 tap assignment 选项；
  \item \cmd{synthesize\_multisource\_clock\_taps}（合并等价时钟单元、将端点赋给最近 tap、沿路径分裂、复制 UPF/SDC/属性）；
  \item \cmd{clock\_opt -from build\_clock -to route\_clock} 构建由 tap 驱动的本地子树，并修复 ignore 例外之外的 DRC；
  \item 分析时钟树结果。
\end{enumerate}

\subsection{使用集成 Tap Assignment 的常规多源时钟树}
\subsubsection*{Implementing a Regular Multisource Clock Tree Using Integrated Tap Assignment}

若在 \cmd{place\_opt} 前插入 tap drivers，工具可在布局优化中分配 sink 以改善总体时序与时钟门控路径。流程与常规类似，但在 mesh 分析后：将 \appopt{place\_opt.flow.enable\_multisource\_clock\_trees} 设为 \cmd{true} 并运行 \cmd{place\_opt}，再 \cmd{clock\_opt -from build\_clock -to route\_clock}。

\subsection{仅含 H-Tree 全局结构的常规多源时钟树}
\subsubsection*{Implementing a Regular Multisource Clock Tree With an H-Tree-Only Global Clock Tree Structure}

全局结构无 mesh，H-tree 直接驱动 tap（图~\ref{fig:mscts-htree-only}）。用 \cmd{set\_regular\_multisource\_clock\_tree\_options} 与 \cmd{synthesize\_regular\_multisource\_clock\_trees} 完成 tap 插入与 H-tree 综合，再做 tap assignment 与 \cmd{clock\_opt}。

\begin{figure}[htbp]
\centering
\fbox{\parbox{0.85\textwidth}{\centering\vspace{1.0cm}
\textit{（原书 Figure 49：Regular Multisource Clock Tree With an H-Tree-Only Global Clock Structure）}
\vspace{1.0cm}}}
\caption{仅含 H-Tree 全局结构的常规多源时钟树}
\label{fig:mscts-htree-only}
\end{figure}

\subsection{实现 MSCTS 多重 H-Tree}
\subsubsection*{Implementing MSCTS Multiple H-Tree}

用 \cmd{set\_regular\_multisource\_clock\_tree\_options} 的 \opt{-htree\_sections} 在 floorplan 不同区域构建多个全局树（图~\ref{fig:mscts-multi-h}），可用对称 tap 配置或坐标指定 tap 位置。

\begin{figure}[htbp]
\centering
\fbox{\parbox{0.85\textwidth}{\centering\vspace{1.0cm}
\textit{（原书 Figure 50/51：MSCTS Multiple H-Tree）}
\vspace{1.0cm}}}
\caption{MSCTS 多重 H-Tree}
\label{fig:mscts-multi-h}
\end{figure}

\subsection{实现功耗感知柔性 H-Tree}
\subsubsection*{Implementing Power Aware Flexible H-Tree}

在柔性 H-Tree 综合中加入功耗感知以优化总时钟功耗与总体 latency。须至少有一个动态功耗 scenario 在 H-tree 综合时为 active；多 scenario 时取最差 active scenario 做功耗代价。初始 tap 配置基于 floorplan、sink 分布、布局与布线约束等估计各子树 latency/功耗并比较。启用：\appopt{cts.multisource.power\_aware\_flexible\_htree\_synthesis}=\cmd{true}。

\subsection{实现结构型多源时钟树}
\subsubsection*{Implementing a Structural Multisource Clock Tree}

子树由 mesh 或预定义驱动直接驱动；通过合并、分裂、改尺寸与搬移优化，同时保留用户定义的级数结构。主要步骤：MV/EM 设置 → \cmd{create\_clock\_straps} → 插入 mesh drivers → 综合全局树 → 用 \cmd{set\_annotated\_delay}/\cmd{set\_annotated\_transition} 建模 mesh 延迟（避免子树综合时看到过大延迟）→ \cmd{synthesize\_multisource\_clock\_subtrees} → \cmd{route\_clock\_straps} → \cmd{analyze\_subcircuit} → \cmd{clock\_opt -from build\_clock -to route\_clock}。

\subsection{使用集成子树综合的 SMSCTS}
\subsubsection*{Implementing a Structural Multisource Clock Tree Using Integrated Subtree Synthesis}

用 \cmd{place\_opt} 综合优化本地子树可改善时序 QoR。额外注意：可启用 \appopt{cts.multisource.enable\_multi\_corner\_support}；用 \cmd{set\_multisource\_clock\_subtree\_options} 指定 \opt{-clock}、\opt{-driver\_objects} 及可选的 \opt{-dont\_merge\_cells}、\opt{-balance\_levels}/\opt{-target\_level}、\opt{-max\_total\_wire\_delay}、\opt{-enable\_icg\_reordering} 等；启用 \appopt{place\_opt.flow.enable\_multisource\_clock\_trees} 后运行 \cmd{place\_opt}；布线到 strap 并分析 mesh 后，启用 \appopt{clock\_opt.flow.enable\_multisource\_clock\_trees} 再跑 \cmd{clock\_opt} 以增量优化子树。

\subsection{改善 SMSCTS 与 CCD 的相关性}
\subsubsection*{Improving the Correlation Between the SMSCTS and CCD}

将 \appopt{cts.multisource.enable\_subtree\_synthesis\_aware\_ccd} 设为 \cmd{true}，可在 \cmd{place\_opt} 与 \cmd{clock\_opt} 的 \cmd{build\_clock} 阶段改善 CCD balance point 实现与 SMSCTS 后时序，降低长短路径违例。

\subsection{插入时钟驱动器}
\subsubsection*{Inserting Clock Drivers}

插入前：指定全部 CTS 设置（含布线规则）；可将 \appopt{cts.multisource.enable\_clock\_driver\_snapping} 设为 \cmd{true} 使驱动与 mesh strap 对齐。\cmd{create\_clock\_drivers} 可用于 mesh drivers（常规/结构型）与 tap drivers（常规）。要点：
\begin{itemize}
  \item \opt{-loads}：负载网或同网引脚/端口；
  \item \opt{-location} 精确位置，或 \opt{-boxes} 网格（可配 \opt{-boundary}）；
  \item \opt{-lib\_cells} 或 \opt{-template}；
  \item mesh drivers 用 \opt{-short\_outputs} 将输出短接驱动 mesh；
  \item \opt{-configuration} 可插入多层驱动形成 H-tree（图~\ref{fig:mscts-drivers}）。
\end{itemize}

\begin{figure}[htbp]
\centering
\fbox{\parbox{0.85\textwidth}{\centering\vspace{1.0cm}
\textit{（原书 Figure 52：Clock Drivers Placed in an H-Tree Structure）}
\vspace{1.0cm}}}
\caption{按 H-Tree 结构放置的时钟驱动器}
\label{fig:mscts-drivers}
\end{figure}

插入后驱动标记为 fixed 与 don't touch；避免与 fixed 单元、blockage、宏重叠，但可与非 fixed 单元重叠。命令不做合法化，需 \cmd{legalize\_placement}（多次插入后统一合法化以节省时间）。布线 H-tree 时使用时钟网规则允许的最高层；为防驱动置于预布线下方导致引脚可达性问题，会为同层预布线创建临时 placement blockage，布线后删除。多行高单元放置困难时可将 \appopt{cts.multisource.enable\_pin\_accessibility\_for\_global\_clock\_trees} 设为 \cmd{false}。移除用 \cmd{remove\_clock\_drivers}。

多级物理层次：用 \cmd{set\_editability} 启用下级编辑；设 \appopt{cts.multisource.enable\_mlph\_flow}=\cmd{true}；插入驱动；用 \cmd{synthesize\_multisource\_global\_clock\_trees -roots -leaves -use\_zroute\_for\_pin\_connections} 布线。下级优先复用既有时钟端口，否则在允许时创建新端口（\cmd{set\_freeze\_ports}）。

\subsection{综合全局时钟树}
\subsubsection*{Synthesizing the Global Clock Trees}

\cmd{synthesize\_multisource\_global\_clock\_trees} 综合并详细布线 H-tree 风格全局树，尽量最小化端点 skew。须指定 \opt{-nets} 与 \opt{-lib\_cells}。默认用 Custom Router；引脚连接 DRC 难解时用 \opt{-use\_zroute\_for\_pin\_connections}；\opt{-skip\_pin\_connections} 可在靠近引脚的最高可用金属层停止。也可仅对已综合结构做详细布线（\opt{-roots}/\opt{-leaves}）。移除用 \cmd{remove\_multisource\_global\_clock\_trees}。

\subsection{创建时钟 Strap}
\subsubsection*{Creating Clock Straps}

\cmd{create\_clock\_straps} 创建单层直线金属形状，可实现：
\begin{itemize}
  \item \textbf{Clock mesh}：水平与垂直层上的二维网格，交叉点用 via 连接（图~\ref{fig:cts-mesh}）；
  \item \textbf{Clock spine}：一维或二维；二维由一维 spine 连接正交方向多条 stripe，不同 spine 的 stripe 不相连，最小间距称 backoff（图~\ref{fig:cts-spine}）。
\end{itemize}

\begin{figure}[htbp]
\centering
\fbox{\parbox{0.85\textwidth}{\centering\vspace{1.0cm}
\textit{（原书 Figure 53：Clock Mesh Structure）}
\vspace{1.0cm}}}
\caption{Clock Mesh 结构}
\label{fig:cts-mesh}
\end{figure}

\begin{figure}[htbp]
\centering
\fbox{\parbox{0.85\textwidth}{\centering\vspace{1.0cm}
\textit{（原书 Figure 54：Clock Spine Structure）}
\vspace{1.0cm}}}
\caption{Clock Spine 结构}
\label{fig:cts-spine}
\end{figure}

主要选项：\opt{-net}、\opt{-boundary}、\opt{-keepouts}（\opt{-allow\_splitting false} 可禁止跨越 keepout 时分裂）、\opt{-layers}、\opt{-width}、\opt{-grid}（单条距离或 start/end/step 迭代列表）、\opt{-margins}、\opt{-create\_ends}、\opt{-allow\_floating}、\opt{-type}（\cmd{user\_route}/\cmd{stripe}/\cmd{detect}）、spine 的 \opt{-length}/\opt{-backoff}/\opt{-spine\_direction}、屏蔽相关 \opt{-bias}/\opt{-bias\_to\_nets}/\opt{-bias\_margins}、\opt{-clear}。

Mesh 示例：
\begin{lstlisting}
icc2_shell> create_clock_straps -nets [get_nets clk1_mesh] \
   -layers {M7 M8} -widths {2.4 2.4} -types {stripe stripe} \
   -grids {{20 1200 100} {60 980 150}} -boundary {{0 0} {1200 980}}
\end{lstlisting}

\subsection{布线到时钟 Strap}
\subsubsection*{Routing to Clock Straps}

\cmd{route\_clock\_straps -nets} 将驱动与负载连到 \cmd{shape\_use}=\cmd{stripe} 的 strap（不连 \cmd{user\_route}）。\opt{-topology}：
\begin{itemize}
  \item \cmd{fishbone}（默认）：各驱动单独连最近 stripe；多负载用 comb 连到 finger 再连 stripe（图~\ref{fig:fishbone}）。可设 \opt{-fishbone\_fanout}/\opt{-fishbone\_span}/\opt{-fishbone\_sub\_span}/\opt{-fishbone\_layers}（图~\ref{fig:fish-fanout}、\ref{fig:fish-span}）。
  \item \cmd{comb}：每引脚直接连最近 stripe；曼哈顿距离超过 comb distance（默认两全局布线格，可用 \appopt{route.common.comb\_distance}）时用 Steiner 连最近 net shape（图~\ref{fig:comb-topo}）。适合大量引脚直接位于 stripe 下，但堆叠 via 可能引起物理 DRC（图~\ref{fig:comb-via}）。
  \item \cmd{sub\_strap}：在中间层创建平行于 stripe 的附加 strap，减少堆叠 via（图~\ref{fig:substrap}）。须 \opt{-sub\_strap\_layers}；可选 \opt{-sub\_strap\_max\_delay}（默认 2~ps）。
\end{itemize}

\begin{figure}[htbp]
\centering
\fbox{\parbox{0.85\textwidth}{\centering\vspace{0.8cm}\textit{（原书 Figure 55：Fishbone Topology）}\vspace{0.8cm}}}
\caption{Fishbone 拓扑}
\label{fig:fishbone}
\end{figure}

\begin{figure}[htbp]
\centering
\fbox{\parbox{0.85\textwidth}{\centering\vspace{0.8cm}\textit{（原书 Figure 56：Fanout of a Finger）}\vspace{0.8cm}}}
\caption{Fishbone finger 的扇出}
\label{fig:fish-fanout}
\end{figure}

\begin{figure}[htbp]
\centering
\fbox{\parbox{0.85\textwidth}{\centering\vspace{0.8cm}\textit{（原书 Figure 57：Span and Subspan）}\vspace{0.8cm}}}
\caption{Fishbone 的 span 与 subspan}
\label{fig:fish-span}
\end{figure}

\begin{figure}[htbp]
\centering
\fbox{\parbox{0.85\textwidth}{\centering\vspace{0.8cm}\textit{（原书 Figure 58：Comb Topology）}\vspace{0.8cm}}}
\caption{Comb 拓扑}
\label{fig:comb-topo}
\end{figure}

\begin{figure}[htbp]
\centering
\fbox{\parbox{0.85\textwidth}{\centering\vspace{0.8cm}\textit{（原书 Figure 59：Vias in the Comb Topology）}\vspace{0.8cm}}}
\caption{Comb 拓扑中的 Via}
\label{fig:comb-via}
\end{figure}

\begin{figure}[htbp]
\centering
\fbox{\parbox{0.85\textwidth}{\centering\vspace{0.8cm}\textit{（原书 Figure 60：Vias in the Substrap Topology）}\vspace{0.8cm}}}
\caption{Substrap 拓扑中的 Via}
\label{fig:substrap}
\end{figure}

\subsection{分析 Clock Mesh}
\subsubsection*{Analyzing the Clock Mesh}

Mesh 需要比传统结构更高的时序精度。用 \cmd{analyze\_subcircuit} 做晶体管级仿真并回标时序。前提：mesh 已详细布线；有门级电路模型与晶体管模型；可访问 NanoSim/FineSim/HSPICE。须指定 \opt{-net}（或 \opt{-to}）与 \opt{-name}；多时钟时用 \opt{-clock}。可选在端口/起点标注转换（时钟须为 propagated）。步骤：RC 提取 → 生成 SPICE → 仿真 → 生成含 \cmd{set\_disable\_timing}/\cmd{set\_annotated\_delay}/\cmd{set\_annotated\_transition} 的注释文件（多驱动 mesh 禁用除 anchor driver 外的驱动）→ 应用注释。各步可用独立选项单独运行。

\subsection{自动 Tap 插入与 H-Tree 综合}
\subsubsection*{Performing Automated Tap Insertion and H-Tree Synthesis}

\begin{enumerate}
  \item 用 \cmd{set\_regular\_multisource\_clock\_tree\_options} 指定 \opt{-clock}、\opt{-topology htree\_only}、\opt{-tap\_boxes}、\opt{-tap\_lib\_cells}、\opt{-htree\_lib\_cells} 等（库单元须已 \cmd{set\_lib\_cell\_purpose -include cts}）。可报告/移除对应 options。
  \item （可选）启用柔性 H-tree：\appopt{cts.multisource.flexible\_htree\_synthesis}=\cmd{true}；对称 tap：\appopt{cts.multisource.tap\_selection}=\cmd{symmetric}。
  \item \cmd{synthesize\_regular\_multisource\_clock\_trees}（阶段 \cmd{tap\_synthesis} 与 \cmd{htree\_synthesis}；可用 \opt{-from}/\opt{-to}）。
\end{enumerate}

\subsection{指定 Tap Assignment 选项与设置}
\subsubsection*{Specifying Tap Assignment Options and Settings}

运行 \cmd{synthesize\_multisource\_clock\_taps} 前：
\begin{enumerate}
  \item 移除不必要的 dont\_touch/size\_only/fixed/locked，以免阻碍合并/分裂；
  \item 用 \cmd{set\_multisource\_clock\_tap\_options} 指定 \opt{-clock}、\opt{-driver\_objects}、\opt{-num\_taps}（须等于驱动数）等；
  \item （可选）用 \cmd{create\_multisource\_clock\_sink\_group} 保持 skew group 对象在同一 tap 下；
  \item （可选）设置合并/分裂命名相关的 \appopt{cts.multisource.subtree\_*} 前缀/后缀。
\end{enumerate}

\subsection{构建本地时钟子树结构}
\subsubsection*{Building the Local Clock Subtree Structures}

结构型多源树中，用 \cmd{synthesize\_multisource\_clock\_subtrees} 综合由 mesh 直接驱动的本地子树。前置：建好全局结构；用真实 annotated delay/transition 建模 mesh；清除不必要的保护属性；用 \cmd{set\_multisource\_clock\_subtree\_options} 配置。可选：多角优化、忽略驱动上 max\_cap/max\_fanout（\appopt{cts.multisource.ignore\_drc\_on\_subtree\_driver}）、fishbone 子树布线（\appopt{cts.multisource.subtree\_routing\_mode}）等。

命令步骤：合并等价单元 → 分裂/改尺寸/搬移并 snap 第一级到 strap → 布线（默认可全局布线；启用 fishbone 则混合详细 fishbone 与全局 comb，见图~\ref{fig:fish-detail}）→ 基于布线细化。可用 \opt{-from}/\opt{-to} 运行 \cmd{merge}/\cmd{optimize}/\cmd{route\_clock}/\cmd{refine}。由此构建的结构型子树在后续 CTS 命令中不会被更改或移除。

\begin{figure}[htbp]
\centering
\fbox{\parbox{0.85\textwidth}{\centering\vspace{1.0cm}
\textit{（原书 Figure 61：Detail Fishbone Routes and Global Comb Routes）}
\vspace{1.0cm}}}
\caption{详细 Fishbone 布线与全局 Comb 布线}
\label{fig:fish-detail}
\end{figure}

% ============================================================
\section{分析时钟树结果}
\subsection*{Analyzing the Clock Tree Results}

综合后分析是否满足要求，通常包括：生成时钟树 QoR 报告；分析时钟时序；在 GUI 中验证时钟实例布局。

\subsection{生成时钟树 QoR 报告}
\subsubsection*{Generating Clock Tree QoR Reports}

用 \cmd{report\_clock\_qor}，可生成：
\begin{itemize}
  \item \textbf{Summary}（默认）：latency、skew、DRC、面积、缓冲计数；
  \item \textbf{Latency}（\opt{-type latency}）：每时钟最长/最短路径；
  \item \textbf{DRC violators}（\opt{-type drc\_violators}）；
  \item \textbf{Robustness}（\opt{-type robustness}，须 \opt{-robustness\_corner}）：当前 corner 与指定 robustness corner 的 latency 比；
  \item \textbf{Balance groups}（\opt{-type balance\_groups}）：含指定与实际 latency offset；
  \item \textbf{Local skew}（\opt{-type local\_skew}）：每时钟/skew group 汇总及最大/最小局部 skew；
  \item \textbf{Power}（\opt{-type power}）：按时钟/scenario 的漏电、内部、sink、网络翻转、动态与总功耗。
\end{itemize}
可用 \opt{-csv}/\opt{-output} 输出 CSV。直方图类型：\opt{-histogram\_type} 为 \cmd{latency}/\cmd{transition}/\cmd{capacitance}/\cmd{local\_skew}/\cmd{robustness}/\cmd{wire\_delay\_fraction}。可用 \opt{-from}/\opt{-to}/\opt{-through}、\opt{-clock}、\opt{-skew\_group}、\opt{-mode}、\opt{-corner}、\opt{-scenario} 限定范围（\opt{-scenario} 与 \opt{-mode}/\opt{-corner} 互斥）。

\subsection{报告时钟树功耗}
\subsubsection*{Reporting Clock Tree Power}

用 \cmd{report\_clock\_power}。可用 \opt{-clocks}/\opt{-modes}/\opt{-corners}/\opt{-scenario} 限定；\opt{-type per\_segment} 或 \opt{-type per\_subtree}：
\begin{itemize}
  \item \textbf{Segment}：从一个 ICG 输出到其扇出中下一 ICG 或 sink 的缓冲树；
  \item \textbf{Subtree}：从 ICG 一直到其扇出全部 sink 的时钟树（图~\ref{fig:seg-sub}）。
\end{itemize}

\begin{figure}[htbp]
\centering
\fbox{\parbox{0.85\textwidth}{\centering\vspace{1.0cm}
\textit{（原书 Figure 62：Segments and Subtrees of a Clock Tree）}
\vspace{1.0cm}}}
\caption{时钟树的 Segment 与 Subtree}
\label{fig:seg-sub}
\end{figure}

\subsection{创建时钟网络引脚集合}
\subsubsection*{Creating Collections of Clock Network Pins}

\cmd{get\_clock\_tree\_pins} 默认返回所有 scenario 全部时钟网上的引脚。可用 \opt{-clocks}、\opt{-scenarios}/\opt{-modes}/\opt{-corners}、\opt{-from}/\opt{-through}/\opt{-to}、\opt{-filter} 限定。例如：
\begin{lstlisting}
icc2_shell> set clk1_icg_pin \
   [get_clock_tree_pins -filter is_on_ICG -clocks clk1]
\end{lstlisting}

\subsection{分析时钟时序}
\subsubsection*{Analyzing Clock Timing}

用 \cmd{report\_clock\_timing}，须指定 \opt{-type}：
\begin{itemize}
  \item \cmd{skew}：单时钟局部 skew；
  \item \cmd{interclock\_skew}：时钟间 skew；
  \item \cmd{latency}：延迟；
  \item \cmd{transition}：转换时间。
\end{itemize}

\subsection{在 GUI 中分析时钟树}
\subsubsection*{Analyzing Clock Trees in the GUI}

IC Compiler II GUI 提供 Clock Tree Visual Mode，可将时钟树信息叠加到版图视图，或在原理图视图中查看结构。详见 \textit{IC Compiler II Graphical User Interface User Guide} 中的 Using Map and Visual Modes。
